%&preamble
\documentclass[11pt]{article}

% Packages for the final build.
\usepackage[T1]{fontenc}
\usepackage{lmodern}
\usepackage{amsmath}
\usepackage{xcolor}
\usepackage[backend=biber,style=authoryear]{biblatex}
\usepackage{tikz}
\usepackage[hidelinks]{hyperref}
\usepackage{cleveref}
% Page and paragraph layout for the final build.
\usepackage[margin=1in]{geometry}
\usepackage{parskip}
\overfullrule=0pt
% Reusable commands and environments for the final build.
\newcommand{\important}[1]{%
	\par\noindent\textcolor{blue!70!black}{\textbf{#1}}\par
}
% Document metadata for the final build.
\title{My LaTeX Document}
\author{Your Name}
\usepackage[en-US]{datetime2}

\newcommand{\CreateDate}{\DTMDate{2026-10-01}}

\date{
    Created on : \CreateDate \\
    Last Modified: \today}

\addbibresource{references.bib}
\csname endofdump\endcsname % <--- Everything ABOVE this line is precompiled into a .fmt file
\includeonly{} % Specify which files to include for faster compilation,if no file is specified, all files are excluded.

\begin{document}
\maketitle

\section{Why modularize?}\label{sec:why-modularize}
Keeping the document setup in small files makes it easier to find and reuse
settings. The main file stays focused on the document's structure and content.

\section{A small example}
Here is a custom command defined in the selected mode's \texttt{commands.tex}:
\important{This text is highlighted using a reusable command.}

Packages loaded in the selected mode's \texttt{packages.tex} also let us typeset equations:
\begin{equation}
	e^{i\pi} + 1 = 0.
\end{equation}

\section{Citation and cross-reference test}\label{sec:reference-test}
\textcite{knuth1984} introduced a foundational work on \TeX. A parenthetical
citation with a page reference looks like this: \parencite[see][chap.~5]{knuth1984}.
The sample \href{https://www.latex-project.org/}{LaTeX Project site} also
provides documentation. A local link jumps to the \hyperref[tab:sample-data]{sample
table}, and \autoref{fig:workflow} exercises automatic figure references.

The data are summarized in \Cref{tab:sample-data}, and the process is shown
in \cref{fig:workflow}. The table and figure can also be referenced together
with \cref{tab:sample-data,fig:workflow}. This sample spans
\crefrange{sec:why-modularize}{sec:reference-test}.

\section{Compile-time benchmark}
These are clean-build wall times in seconds, measured with separate job names;
the final build includes Biber and TikZ rendering. In these runs, final mode's
median was about 3.39 times the draft median.

\begin{table}[htbp]
\centering
\caption{Three clean-build timings per mode, in seconds.}
\label{tab:compile-benchmark}
\begin{tabular}{lrrrr}
\hline
Mode & Run 1 & Run 2 & Run 3 & Median \\
\hline
Draft & 16.37 & 10.70 & 13.04 & 13.04 \\
Final & 44.14 & 38.68 & 57.75 & 44.14 \\
\hline
\end{tabular}
\end{table}

\begin{table}[htbp]
\centering
\caption{Sample measurements used for the compile test.}
\label{tab:sample-data}
\begin{tabular}{lrr}
\hline
Mode & Citations & References \\
\hline
Draft & 2 & 2 \\
Final & 2 & 2 \\
\hline
\end{tabular}
\end{table}

\begin{figure}[htbp]
\centering
\begin{tikzpicture}
\node[draw] (source) at (0,0) {Source};
\node[draw] (build) at (3,0) {Build};
\node[draw] (output) at (6,0) {PDF};
\draw[->] (source) -- (build);
\draw[->] (build) -- (output);
\end{tikzpicture}
\caption{A simple workflow diagram for the figure-reference test.}
\label{fig:workflow}
\end{figure}

Each of \texttt{preamble/draft/} and \texttt{preamble/final/} contains:
\begin{itemize}
	\item \texttt{packages.tex}: package imports
	\item \texttt{layout.tex}: page and paragraph layout
	\item \texttt{commands.tex}: custom commands and environments
	\item \texttt{metadata.tex}: title, author, and date
\end{itemize}

\begin{equation}
    \label{eq:ns}
    \nabla \cdot \mathbf{v} = 0
\end{equation}

\begin{equation}
    \label{eq:test}
    \begin{aligned}
    \frac{\partial{f}}{\partial{x}}+1=0 \\
    f(0) = 0
    \end{aligned}
\end{equation}
The rest of the text
\printbibliography

\end{document}